\chapter{字符串算法}\label{ux5b57ux7b26ux4e32ux7b97ux6cd5}

\section{选择算法}\label{ux9009ux62e9ux7b97ux6cd5}

\begin{longtable}[]{@{}lll@{}}
\toprule\noalign{}
需求 & 算法 & 复杂度 \\
\midrule\noalign{}
\endhead
\bottomrule\noalign{}
\endlastfoot
单模式匹配 & KMP / Z & \passthrough{\lstinline!O(n+m)!} \\
多模式匹配 & AC 自动机 & \passthrough{\lstinline!O(text + 总模式长)!} \\
子串相等、二分答案 & 双哈希 & 期望 \passthrough{\lstinline!O(1)!} 查询 \\
最长回文子串 & Manacher & \passthrough{\lstinline!O(n)!} \\
所有后缀排序/不同子串 & 后缀数组 + LCP & \passthrough{\lstinline!O(n log n)!} \\
子串出现次数/自动机 & SAM & \passthrough{\lstinline!O(n)!} \\
\end{longtable}

\section{Knuth--Morris--Pratt 字符串匹配算法（KMP）}\label{knuthmorrispratt-ux5b57ux7b26ux4e32ux5339ux914dux7b97ux6cd5kmp}

\passthrough{\lstinline!nxt[i]!} 表示模式前缀 \passthrough{\lstinline!p[0..i]!} 的最长真前后缀长度。匹配失败时跳到 \passthrough{\lstinline!nxt[j-1]!}，不回退文本指针；可求出现位置、最小周期、前缀函数树。

最小周期：若 \passthrough{\lstinline!n \% (n - pi[n-1]) == 0!}，周期为 \passthrough{\lstinline!n-pi[n-1]!}，否则为 \passthrough{\lstinline!n!}。

\begin{lstlisting}[language={C++}]
// 接口：prefix_function(s)：返回 KMP 前缀函数 pi；pi[i] 是 s[0..i] 的最长真 border 长度；返回类型 vector<int>。
// 参数：源点编号或输入字符串 s（见本节类型）。
vector<int> prefix_function(const string& s) {
    // 变量：pi=KMP 前缀函数 pi。
    vector<int> pi(s.size());
    // 变量：i=当前循环下标 i。
    for (int i = 1; i < (int)s.size(); ++i) {
        // 变量：j=内层循环下标 j。
        int j = pi[i - 1];
        while (j && s[i] != s[j]) j = pi[j - 1];
        if (s[i] == s[j]) ++j;
        pi[i] = j;
    }
    return pi;
}
// 接口：kmp(text, pat)：返回模式串 pat 在 text 中所有 0-based 起点；返回类型 vector<int>。
// 参数：主串 text；模式串 pat。
vector<int> kmp(const string& text, const string& pat) {
    // 变量：pi=模式串 pat 的 KMP 前缀函数数组 pi；pos=所有匹配起点的返回数组 pos。
    vector<int> pi = prefix_function(pat), pos;
    if (pat.empty()) return pos; // 本接口约定空模式不返回匹配位置。
    // 变量：j=当前已经匹配的模式串前缀长度 j。
    int j = 0;
    // 变量：i=当前循环下标 i。
    for (int i = 0; i < (int)text.size(); ++i) {
        while (j && text[i] != pat[j]) j = pi[j - 1];
        if (text[i] == pat[j]) ++j;
        if (j == (int)pat.size()) {
            pos.push_back(i - (int)pat.size() + 1);
            j = pi[j - 1];
        }
    }
    return pos;
}
\end{lstlisting}

\section{Z 函数}\label{z-ux51fdux6570}

\passthrough{\lstinline!z[i]!} 是 \passthrough{\lstinline!s[i..]!} 与 \passthrough{\lstinline!s[0..]!} 的最长公共前缀，本模板约定 \passthrough{\lstinline!z[0]=n!}。维护当前最右匹配段 \passthrough{\lstinline![l,r)!}，落在段内时先继承 \passthrough{\lstinline!z[i-l]!}。用 Z 函数做匹配时，分隔符必须保证不在输入字符集中。

\begin{lstlisting}[language={C++}]
// 接口：z_function(s)：返回 Z 数组；z[i] 是 s 与 s[i..] 的最长公共前缀长度；返回类型 vector<int>。
// 参数：源点编号或输入字符串 s（见本节类型）。
vector<int> z_function(const string& s) {
    // 变量：n=当前元素/顶点/状态数量 n；l=当前区间左端点 l；r=当前区间右端点 r；z=Z 函数数组 z；z[i] 是 s 与 s[i..] 的最长公共前缀。
    int n = s.size(), l = 0, r = 0; vector<int> z(n);
    if (n) z[0] = n;
    // 变量：i=当前循环下标 i。
    for (int i = 1; i < n; ++i) {
        if (i < r) z[i] = min(r - i, z[i - l]);
        while (i + z[i] < n && s[z[i]] == s[i + z[i]]) ++z[i];
        if (i + z[i] > r) l = i, r = i + z[i];
    }
    return z;
}
\end{lstlisting}

\section{字符串哈希}\label{ux5b57ux7b26ux4e32ux54c8ux5e0c}

前缀哈希 \passthrough{\lstinline!H[i+1]=H[i]*B+s[i]!}，子串哈希：\passthrough{\lstinline!H[r]-H[l]*powB[r-l]!}。双模数或 64 位自然溢出降低碰撞；随机底数并不能提供形式化保证。比较子串、二分最长公共前缀、回文判定都可用。

\begin{lstlisting}[language={C++}]
// 变量：BASE=字符串哈希使用的固定底数 BASE。
const unsigned long long BASE = 911382323ull;
// 变量：h=“字符串哈希”这段代码中的临时量 h；值由本行初始化式确定；pw=幂次预处理数组 pw。
vector<unsigned long long> h(n + 1), pw(n + 1, 1);
// 变量：i=当前循环下标 i。
for (int i = 0; i < n; ++i)
    h[i + 1] = h[i] * BASE + (unsigned char)s[i], pw[i + 1] = pw[i] * BASE;
// 接口：get_hash(l, r)：供“字符串哈希”当前语句回调；返回值含义见调用处条件。
// 参数：闭区间左端点 l；闭区间右端点 r。
auto get_hash = [&](int l, int r) { return h[r] - h[l] * pw[r - l]; }; // [l,r)
\end{lstlisting}

\section{Manacher}\label{manacher}

把每个字节转成非负整数，在字符之间插入 \passthrough{\lstinline!-1!}，并在两端放不同的负数哨兵，求每个中心向外的回文半径 \passthrough{\lstinline!p[i]!}；维护最右回文段 \passthrough{\lstinline!(center,right)!}，线性求最长回文。整数哨兵不会与任意输入字节冲突。

\begin{lstlisting}[language={C++}]
// 变量：t=插入分隔符和两端哨兵后的整数串 t。
vector<int> t = {-2, -1};
// 变量：c=当前输入字节 c。
for (unsigned char c : s) t.push_back(c), t.push_back(-1);
t.push_back(-3);
// 变量：p=Manacher 半径数组 p；center=Manacher 当前最右回文的中心 center；right=滑动窗口右端点 right。
vector<int> p(t.size()); int center = 0, right = 0;
// 变量：i=当前循环下标 i。
for (int i = 1; i + 1 < (int)t.size(); ++i) {
    // 变量：mir=Manacher 中与 i 关于中心对称的位置 mir。
    int mir = 2 * center - i;
    if (i < right) p[i] = min(right - i, p[mir]);
    while (t[i + 1 + p[i]] == t[i - 1 - p[i]]) ++p[i];
    if (i + p[i] > right) center = i, right = i + p[i];
}
\end{lstlisting}

\section{字典树与 Aho--Corasick 多模式匹配自动机（AC 自动机）}\label{ux5b57ux5178ux6811ux4e0e-ahocorasick-ux591aux6a21ux5f0fux5339ux914dux81eaux52a8ux673aac-ux81eaux52a8ux673a}

Trie 每个节点保存儿子、终止次数。AC 自动机用 BFS 建 fail：节点的 fail 指向最长真后缀状态，匹配时沿 fail 转移；把输出计数沿 fail 累加可统计所有模式出现次数。模式集合固定后一次建树、多次查询。

\begin{lstlisting}[language={C++}]
struct AC {
    int ch[MAXN][26]{}; // Trie 边；build 后也会补齐为自动机转移。
    int fail[MAXN]{};   // 最长真后缀节点。
    int out[MAXN]{};    // 模式串结束次数及 fail 链汇总次数。
    int tot = 0;          // 当前最后一个节点编号；根为 0，节点范围 [0,tot]。
    // 接口：insert(s)：把参数表示的数/字符串/版本插入当前结构；会修改当前结构；多测时重新构造或显式清空成员。
    // 参数：源点编号或输入字符串 s（见本节类型）。
    void insert(const string& s) {
        // 变量：u=当前顶点/状态编号 u。
        int u = 0;
        // 变量：c=当前字符、容量或第三个操作数 c。
        for (char c : s) {
            // 变量：x=当前输入值/查询值/坐标 x。
            int x = c - 'a';
            if (!ch[u][x]) ch[u][x] = ++tot;
            u = ch[u][x];
        }
        // 重复模式串要累加，而不是只记一个 bool。
        ++out[u];
    }
    // 接口：build()：由当前输入/成员状态完成数据结构预处理；完成初始化；多测时每组重新调用。
    void build() {
        // 变量：q=当前 BFS/单调算法使用的队列 q。
        queue<int> q;
        // 根节点的直接儿子的 fail 都是根。
        // 变量：c=当前字符、容量或第三个操作数 c。
        for (int c = 0; c < 26; ++c) if (ch[0][c]) q.push(ch[0][c]);
        while (!q.empty()) {
            // 变量：u=当前顶点/状态编号 u。
            int u = q.front(); q.pop();
            // 如果要统计每个模式串，推荐建 fail 树后再统一 DFS 汇总；
            // 这里直接累加适合只问“当前位置结束的所有模式数”。
            out[u] += out[fail[u]];
            // 变量：c=当前字符、容量或第三个操作数 c。
            for (int c = 0; c < 26; ++c) {
                // 变量：v=当前相邻顶点/下一状态 v。
                int &v = ch[u][c];
                if (v) {
                    // v 的 fail 是 fail[u] 沿同字符转移。
                    fail[v] = ch[fail[u]][c], q.push(v);
                } else {
                    // 补全转移，扫描时不需要反复 while 回退。
                    v = ch[fail[u]][c];
                }
            }
        }
    }
    // 接口：query(s)：返回本节数据结构在给定位置/区间上的查询结果；返回类型 int。
    // 参数：源点编号或输入字符串 s（见本节类型）。
    int query(const string& s) {
        // 变量：u=当前顶点/状态编号 u；ans=当前累计答案 ans。
        int u = 0, ans = 0;
        // 变量：c=当前字符、容量或第三个操作数 c。
        for (char c : s) {
            // u 表示当前文本前缀的最长 Trie 后缀。
            u = ch[u][c - 'a'];
            ans += out[u];
        }
        return ans;
    }
};
\end{lstlisting}

\section{后缀数组与最长公共前缀（Suffix Array and Longest Common Prefix，LCP）}\label{ux540eux7f00ux6570ux7ec4ux4e0eux6700ux957fux516cux5171ux524dux7f00suffix-array-and-longest-common-prefixlcp}

倍增排序：按 \passthrough{\lstinline!(rank[i], rank[i+k])!} 对后缀排序，重复 \passthrough{\lstinline!O(log n)!} 次，复杂度 \passthrough{\lstinline!O(n log²n)!}（用计数排序可到 \passthrough{\lstinline!O(n log n)!}）。Kasai 求相邻后缀 LCP，RMQ 查询任意两个后缀 LCP。常见应用：不同子串数 \passthrough{\lstinline!n(n+1)/2 - ΣLCP!}、最长重复子串、字典序第 k 个子串。

\begin{lstlisting}[language={C++}]
// 接口：suffix_array(s)：返回 sa；sa[i] 是字典序第 i 小后缀的 0-based 起点；返回类型 vector<int>。
// 参数：源点编号或输入字符串 s（见本节类型）。
vector<int> suffix_array(string s) {
    // 加入最小哨兵，保证所有真正后缀都排在哨兵之后。
    // 变量：n=当前元素/顶点/状态数量 n。
    s.push_back(0); int n = s.size();
    // 变量：sa=后缀数组 sa；sa[i] 是第 i 小后缀起点；r=当前区间右端点 r；nr=本轮扩展得到的新右边界 nr。
    vector<int> sa(n), r(n), nr(n);
    iota(sa.begin(), sa.end(), 0);
    // 变量：i=当前循环下标 i。
    for (int i = 0; i < n; ++i) r[i] = (unsigned char)s[i];
    // 变量：k=当前处理的二进制位/倍增层 k。
    for (int k = 1;; k <<= 1) {
        // 用长度 k 的前半排名和后半排名比较长度 2k 的前缀。
        // 接口：比较器 lambda(i, j)：供“后缀数组与最长公共前缀（Suffix Array and Longest Common Prefix，LCP）”当前语句回调；返回值含义见调用处条件。
        // 参数：当前 0/1-based 位置 i（见接口区间约定）；比较器收到的另一个 0-based 下标 j。
        sort(sa.begin(), sa.end(), [&](int i, int j) {
            if (r[i] != r[j]) return r[i] < r[j];
            // 变量：ri=后缀 i 的当前第一关键字排名 ri。
            int ri = i + k < n ? r[i + k] : -1;
            // 变量：rj=后缀 j 的当前第一关键字排名 rj。
            int rj = j + k < n ? r[j + k] : -1;
            return ri < rj;
        });
        // 重新压缩排名；二元组完全相同的后缀共享排名。
        nr[sa[0]] = 0;
        // 变量：i=当前循环下标 i。
        for (int i = 1; i < n; ++i)
            nr[sa[i]] = nr[sa[i - 1]] +
                (r[sa[i - 1]] != r[sa[i]] ||
                 (sa[i - 1] + k < n ? r[sa[i - 1] + k] : -1) !=
                 (sa[i] + k < n ? r[sa[i] + k] : -1));
        r.swap(nr);
        if (r[sa.back()] == n - 1) break;
    }
    sa.erase(sa.begin()); return sa;
}
\end{lstlisting}

\section{后缀自动机（Suffix Automaton，SAM）}\label{ux540eux7f00ux81eaux52a8ux673asuffix-automatonsam}

在线加入字符，状态保存 \passthrough{\lstinline!len!}、\passthrough{\lstinline!link!}、转移；每个状态代表一组 endpos 相同的子串。状态数最多 \passthrough{\lstinline!2n-1!}。按 \passthrough{\lstinline!len!} 降序累加出现次数，可求每个子串出现频次、不同子串数 \passthrough{\lstinline!Σ(len[v]-len[link[v]])!}。转移字符集大时用 \passthrough{\lstinline!map!}，小字符集用数组。

\begin{lstlisting}[language={C++}]
struct SAM {
    struct Node {
        int next[26]{}; // 从该 endpos 等价类沿字符转移。
        int link = -1;  // suffix link，指向更短的等价类。
        int len = 0;    // 该状态能代表的最长子串长度。
        int occ = 0;    // 直接作为前缀末尾出现的次数，之后沿 link 汇总。
    } st[2 * MAXN];
    int last = 0; // 当前整个字符串对应的状态。
    int sz = 1;   // 状态 0 是初始状态。
    // 接口：extend(c)：向当前字符串自动机加入一个字符并更新状态；会修改当前结构；多测时重新构造或显式清空成员。
    // 参数：字符编号/边容量/候选对象 c（见本节）。
    void extend(char c) {
        // cur 表示加入 c 后的新前缀及其所有后缀。
        // 变量：cur=当前扫描位置的累计状态 cur；p=当前质数、节点编号或指针 p；具体角色见所在公式。
        int cur = sz++, p = last; st[cur].len = st[last].len + 1; st[cur].occ = 1;
        // 从 last 沿 suffix link 补齐缺失的 c 转移。
        while (p >= 0 && !st[p].next[c - 'a']) st[p].next[c - 'a'] = cur, p = st[p].link;
        if (p < 0) st[cur].link = 0;
        else {
            // 变量：q=当前查询对象 q。
            int q = st[p].next[c - 'a'];
            if (st[p].len + 1 == st[q].len) st[cur].link = q;
            else {
                // q 太长，复制出 clone 把长度截断到 st[p].len+1。
                // 变量：clone=后缀自动机为拆分转移而复制的克隆状态 clone。
                int clone = sz++; st[clone] = st[q]; st[clone].len = st[p].len + 1; st[clone].occ = 0;
                // 把原来能经过 c 到 q 的一段 link 链改指 clone。
                while (p >= 0 && st[p].next[c - 'a'] == q) st[p].next[c - 'a'] = clone, p = st[p].link;
                st[q].link = st[cur].link = clone;
            }
        }
        last = cur;
    }
};
\end{lstlisting}

\section{回文自动机（Palindromic Automaton，PAM，了解）}\label{ux56deux6587ux81eaux52a8ux673apalindromic-automatonpamux4e86ux89e3}

每个节点代表一个不同回文串，fail 指向最长回文后缀；在线插入字符即可统计不同回文子串、每个前缀的回文后缀。代码比 Manacher 长，只有在``动态加入/回文后缀结构''时使用。

\begin{lstlisting}[language={C++}]
struct PAM {
    int next[MAXN][26]; // 回文串两端加字符后的转移。
    int fail[MAXN];     // 最长真回文后缀。
    int len[MAXN];      // 节点对应回文串长度。
    int s[MAXN];        // 已插入字符，寻找可扩展回文时要回看它。
    int occ[MAXN];      // 先记作最长回文后缀的次数，再沿 fail 汇总。
    int last, n, tot;   // last=当前最长回文后缀，n=长度，tot=节点数。
    // 接口：init()：清空成员并建立本自动机所需的初始根节点；完成初始化；多测时每组重新调用。
    void init() {
        // 0 号根长度 0，1 号根长度 -1；两个根让边界转移统一。
        memset(next, 0, sizeof(next)); memset(occ, 0, sizeof(occ));
        n = 0; tot = 1; last = 0; s[0] = -1;
        len[0] = 0; len[1] = -1; fail[0] = fail[1] = 1;
    }
    // 接口：get_fail(u)：沿 fail 链找到能被当前字符继续扩展的节点；返回类型 int。
    // 参数：顶点/状态编号 u。
    int get_fail(int u) {
        // 沿 fail 跳，直到 u 的两端可以接上新字符 s[n]。
        while (s[n - len[u] - 1] != s[n]) u = fail[u];
        return u;
    }
    // 接口：add(c)：执行本节定义的单点/区间增加或加入操作；会修改当前结构；多测时重新构造或显式清空成员。
    // 参数：字符编号/边容量/候选对象 c（见本节）。
    void add(char c) {
        // 先把新字符写入串，再寻找能够扩展的最长回文后缀。
        // 变量：u=当前顶点/状态编号 u。
        s[++n] = c - 'a'; int u = get_fail(last);
        if (!next[u][s[n]]) {
            // 这是一个此前未出现过的新回文节点。
            // 变量：v=当前相邻顶点/下一状态 v。
            int v = ++tot; len[v] = len[u] + 2;
            // 新节点的 fail：从 fail[u] 开始寻找也能被该字符扩展的节点。
            fail[v] = next[get_fail(fail[u])][s[n]];
            next[u][s[n]] = v;
        }
        // last 是当前前缀的最长回文后缀；每个位置只直接计数一次。
        last = next[u][s[n]]; ++occ[last];
    }
    // 接口：build(str)：由当前输入/成员状态完成数据结构预处理；完成初始化；多测时每组重新调用。
    // 参数：待完整构建的字符串 str。
    int build(const string& str) {
        init();
        // 变量：c=当前字符、容量或第三个操作数 c。
        for (char c : str) add(c);
        // fail 从长回文指向短回文，必须按长度降序汇总；
        // 许多简化板子用编号倒序，只在编号也随长度单调时才安全。
        // 变量：ord=按长度/拓扑顺序排列的状态编号 ord。
        vector<int> ord(max(0, tot - 1));
        iota(ord.begin(), ord.end(), 2); // 0/1 是两个虚拟根，不参与统计。
        // 接口：比较器 lambda(a, b)：供“回文自动机（Palindromic Automaton，PAM，了解）”当前语句回调；返回值含义见调用处条件。
        // 参数：输入值/数组 a（见本节定义）；输入值/数组 b（见本节定义）。
        sort(ord.begin(), ord.end(), [&](int a, int b) { return len[a] > len[b]; });
        // 变量：v=当前相邻顶点/下一状态 v。
        for (int v : ord) occ[fail[v]] += occ[v];
        return tot - 1; // 不同非空回文子串数量
    }
};
\end{lstlisting}

\section{字符串常见组合}\label{ux5b57ux7b26ux4e32ux5e38ux89c1ux7ec4ux5408}

\begin{itemize}
\tightlist
\item
  子串出现次数：KMP 的前缀函数树或 SAM 的 endpos 计数。
\item
  多个模式在文本中出现：AC 自动机；需要撤销/动态模式时考虑 fail 树 + 树状数组。
\item
  比较循环字符串：最小表示法（Booth）\passthrough{\lstinline!O(n)!}。
\item
  最短覆盖/最小表示：先判周期，再用 KMP/哈希做边界验证。
\end{itemize}

\begin{lstlisting}[language={C++}]
// 接口：min_rotation(s)：返回最小循环表示的 0-based 起点；返回类型 int。
// 参数：源点编号或输入字符串 s（见本节类型）。
int min_rotation(const string& s) {
    // 变量：t=测试用例数量或当前临时值 t；i=当前循环下标 i；j=内层循环下标 j；n=当前元素/顶点/状态数量 n。
    string t = s + s; int i = 0, j = 1, n = s.size();
    while (i < n && j < n) {
        // 变量：k=当前排名、选取数量或循环层数 k。
        int k = 0;
        while (k < n && t[i + k] == t[j + k]) ++k;
        if (k == n) break;
        if (t[i + k] > t[j + k]) i += k + 1, i += (i == j);
        else j += k + 1, j += (i == j);
    }
    return min(i, j);
}
\end{lstlisting}
